\begin{lemma}
\label{lem:p071-finite-exact-star-polyhedral-insertion}
Let \(A\) and \(F\) be disjoint finite sets with \(A\cup F\ne\varnothing\).
Form a compact exact star \(D\) with centre \(\eta\), exceptional branches
\[
                  [0,R_a]\quad(a\in A),
\]
and full branches
\[
                  [0,+\infty]\quad(b\in F),
\]
all parameterized outwards from \(\eta\).  Let \(Q\) be a nonempty compact
convex polytope.

Suppose that \(u:D\times Q\to\mathbf R\cup\{-\infty\}\) is upper
semicontinuous, is jointly extended-real convex on every finite branch
prism, is nonincreasing in the radial variable on every full branch, and
has the exact trace at each full endpoint.  Write
\(c(y)=u(\eta,y)\), and assume the following exact product junction
inequalities.  For every \(S\subseteq F\), every choice of positive radii
\(s_i\) on the branches \(i\in A\cup S\), and points \(y_i\in Q\), put
\[
 \Lambda=\sum_{i\in A\cup S}\frac1{s_i},
 \qquad
 y_0=\Lambda^{-1}\sum_{i\in A\cup S}\frac{y_i}{s_i}.
\]
Whenever \(\Lambda>0\) and \(c(y_0)\) is finite, require
\begin{equation}
 \Lambda c(y_0)\leq
       \sum_{i\in A\cup S}\frac{u_i(s_i,y_i)}{s_i}.
 \label{eq:p071-finite-exact-star-junction}
\end{equation}

If \(\phi\in C(D\times Q)\) and \(u\leq\phi\), then there is a largest
exact product-comparison function \(E_\phi\) below \(\phi\).  It is
real-valued and continuous on \(D\times Q\), including every full endpoint,
and
\begin{equation}
                              u\leq E_\phi\leq\phi.
 \label{eq:p071-finite-exact-star-insertion}
\end{equation}
Moreover, the operator \(\phi\mapsto E_\phi\) is nondecreasing and
nonexpanding in the uniform norm.
\end{lemma}

\begin{proof}
We first give the Jensen families defining the envelope.  A finite atomic
probability measure \(\mu\) will be supported at node points
\((\eta,y)\) and at finite branch points \((i,s,y)\).  Put
\[
 q_i(\mu)=\int_{\{i\}\times(0,\infty)\times Q}s\,\mathrm d\mu,
 \qquad
 \bar y(\mu)=\int y\,\mathrm d\mu.
\]
For a node root \((\eta,y)\), call \(\mu\) admissible if
\(\bar y(\mu)=y\) and there is a number \(\lambda\geq0\) such that
\begin{equation}
 q_a(\mu)=\lambda\quad(a\in A),
 \qquad
 q_b(\mu)\leq\lambda\quad(b\in F).
 \label{eq:p071-finite-exact-node-moments}
\end{equation}
For a finite point \((a_0,s,y)\) on an exceptional branch, replace the
first group of equalities by
\begin{equation}
 q_{a_0}(\mu)=s+\lambda,\qquad
 q_a(\mu)=\lambda\quad(a\in A\setminus\{a_0\}),
 \label{eq:p071-finite-exact-exception-moments}
\end{equation}
and retain the full-branch inequalities in
\eqref{eq:p071-finite-exact-node-moments}.  For a finite point
\((b_0,s,y)\) on a full branch, require
\begin{equation}
 \begin{split}
 q_a(\mu)&=\lambda \quad(a\in A),\\
 q_b(\mu)&\leq\lambda \quad(b\in F\setminus\{b_0\}),\\
 q_{b_0}(\mu)&\leq s+\lambda .
 \end{split}
 \label{eq:p071-finite-exact-full-moments}
\end{equation}
At the endpoint \((b_0,+\infty,y)\), impose the first two lines of
\eqref{eq:p071-finite-exact-full-moments} and no upper bound on
\(q_{b_0}(\mu)\).  Denote the resulting family by
\(\mathcal J_x\), where \(x\) is the root, and define
\begin{equation}
             E_\phi(x)=\inf_{\mu\in\mathcal J_x}\int\phi\,\mathrm d\mu.
 \label{eq:p071-finite-exact-moment-envelope}
\end{equation}
Every finite-root family contains the singleton at its root.  At a full
endpoint, singletons at finite points of that branch are admissible and
approach the endpoint.  Thus the families are nonempty and the envelope is
finite.

We next prove continuity.  For \(N\geq1\), restrict the support of every
measure on a full branch to radii at most \(N\), allow arbitrary Borel
probability measures on the resulting compact finite star, and call the
corresponding value \(E_{\phi,N}\).  This does not change the value if one
restricts again to finite atomic measures.  Indeed, Carath\'eodory's theorem
applied to all radial moments, the affine coordinates on \(Q\), and the
cost replaces any probability measure by a finite one with the same data.

The admissible-measure correspondence for \(E_{\phi,N}\) is compact-valued
and closed.  When \(A\ne\varnothing\), the common number \(\lambda\) is one
of the bounded exceptional radial moments.  When \(A=\varnothing\), the
existential inequalities impose no additional closedness issue.  Therefore
convergent roots and minimizing measures give lower semicontinuity.
Mixtures of admissible measures show joint convexity on every branch prism.
On a full branch, increasing the root radius relaxes only the last
inequality in \eqref{eq:p071-finite-exact-full-moments}; hence
\(E_{\phi,N}\) is nonincreasing there.  It is constant for root radii
\(s\geq N\), including at \(+\infty\).  A finite convex function on a
compact polytope is upper semicontinuous at its boundary by a triangulation
with the target point as a vertex.  Consequently \(E_{\phi,N}\) is
continuous on each branch prism.  The moment conditions at radius zero are
exactly \eqref{eq:p071-finite-exact-node-moments}; since there are only
finitely many branches, the branchwise continuities glue at the node.

Set
\[
 \omega_N=\max_{b\in F}\sup\{
   |\phi_b(s,y)-\phi_b(s',y)|:
   s,s'\in[N,+\infty],\ y\in Q\},
\]
with \(\omega_N=0\) when \(F=\varnothing\).  Continuity on the finitely many
compactified full prisms gives \(\omega_N\to0\).  The restricted measure
family is smaller, so \(E_\phi\leq E_{\phi,N}\).  Conversely, in a finite
family admissible for \(E_\phi(x)\), replace every full-branch radius
\(s>N\) by \(N\), without changing its mass or its \(Q\)-coordinate.
Every full radial moment decreases.  All full moment conditions in
\eqref{eq:p071-finite-exact-node-moments} and
\eqref{eq:p071-finite-exact-full-moments} are upper bounds, so the new
family is still admissible.  Its cost changes by at most \(\omega_N\).
Taking infima gives the uniform estimate
\begin{equation}
             0\leq E_{\phi,N}-E_\phi\leq\omega_N
             \quad\text{on }D\times Q.
 \label{eq:p071-finite-exact-uniform-truncation}
\end{equation}
Thus \(E_\phi\) is a uniform limit of continuous functions and is
continuous on the whole compact product.  The singleton inequalities at
finite points and passage to a full endpoint now give \(E_\phi\leq\phi\).

We verify that \(E_\phi\) satisfies the comparison conditions.  Convexity
on branch prisms follows by mixing root measures and their corresponding
\(\lambda\)'s.  The same mixing and
\eqref{eq:p071-finite-exact-full-moments} show that the envelope is
nonincreasing on a full branch.  Its exact endpoint trace follows from
continuity.

For the junction condition, normalize the weights in
\eqref{eq:p071-finite-exact-star-junction} by putting
\[
                       \alpha_i=\frac1{\Lambda s_i}.
\]
Choose nearly minimizing measures \(\mu_i\) for the branch roots and let
\(\lambda_i\) be their moment parameters.  Their mixture
\(\mu=\sum_i\alpha_i\mu_i\) has the prescribed \(Q\)-barycentre.  Put
\[
                  L=\sum_i\alpha_i\lambda_i,
                  \qquad q_0=\alpha_i s_i=\Lambda^{-1}.
\]
For every exceptional branch its radial moment in \(\mu\) equals
\(L+q_0\).  For an included full branch the moment is at most \(L+q_0\),
and for an omitted full branch it is at most \(L\).  Hence \(\mu\) is
admissible for the node root, with parameter \(L+q_0\).  Formula
\eqref{eq:p071-finite-exact-moment-envelope}, followed by the
near-minimizing limit, gives
\[
                   E_\phi(\eta,y_0)
                      \leq\sum_i\alpha_iE_{\phi,i}(s_i,y_i),
\]
which is precisely the required junction inequality.

It remains to prove that the moment families encode every consequence of
the primitive comparison inequalities.  We prove the following Jensen
claim:
\begin{equation}
             v(x)\leq\int v\,\mathrm d\mu
 \quad\text{for every exact comparison function \(v\leq\phi\) and
 every \(\mu\in\mathcal J_x\).}
 \label{eq:p071-finite-exact-jensen-claim}
\end{equation}
It is enough to treat finite real-valued $v$.  Indeed, if $v\leq\phi$,
choose constants $K_j<\min\phi$ with $K_j\downarrow-\infty$, and put
$v_j=\max\{v,K_j\}$.  Each $v_j$ is finite, remains below $\phi$,
and satisfies all comparison conditions: the scalar maximum map is convex
and nondecreasing, so it preserves every normalized Jensen inequality.
The finite-valued claim gives
\[
                         v_j(x)\leq\int v_j\,\mathrm d\mu.
\]
The measure $\mu$ has finite support.  Hence the left side decreases to
$v(x)$, while the right side decreases to the extended-real finite sum
$\int v\,\mathrm d\mu$, with zero-weight atoms omitted.  Passing to the
limit proves the claim for $v$.

First let the root be a node.  Ignore atoms of \(\mu\) at radius zero.
For each branch, partition an interval of length \(q_i(\mu)\) into pieces
of lengths equal to the atom weights times their radii, and label every
piece by the corresponding atom.  Use \([0,\lambda]\) for each exceptional
branch and place the labelled interval of a full branch inside
\([0,\lambda]\), which is possible by
\eqref{eq:p071-finite-exact-node-moments}.  At almost every
\(\tau\in[0,\lambda]\), select the labelled atom on every exceptional
branch and on each full branch whose labelled set contains \(\tau\).
Writing the selected atom as \((i,s_i(\tau),y_i(\tau))\), apply
\eqref{eq:p071-finite-exact-star-junction} with reciprocal weights
\(1/s_i(\tau)\).  After integration in \(\tau\), every positive-radius
atom of weight \(\beta\) contributes exactly
\[
             \int_{\text{its interval}}\frac{\mathrm d\tau}{s}
                   =\beta
\]
times its value and its \(Q\)-coordinate.  Convexity of the centre slice
combines the integrated centre values, and a second application of that
convexity combines the original radius-zero atoms.  Since
\(\bar y(\mu)=y\), the result is
\[
                         c(y)\leq\int v\,\mathrm d\mu.
\]
This also covers \(A=\varnothing\); values of \(\tau\) at which no branch
is selected contribute zero.

Now let the root lie on a branch \(i_0\) at radius \(s\).  If \(i_0\in A\),
put \(r=s\).  If \(i_0\in F\), put
\[
                         r=\max\{0,q_{i_0}(\mu)-\lambda\}\leq s.
\]
Choose a submeasure \(\nu\) of the \(i_0\)-branch part of \(\mu\) whose
radial moment is \(r\), splitting one atom if necessary.  The residual
measure \(\mu-\nu\), after normalization when it has positive mass,
satisfies the node moment conditions.  If
\(\beta=\nu(D\times Q)\) and both \(\beta\) and \(1-\beta\) are positive,
branchwise Jensen and the node case give
\[
 v_{i_0}\left(r,\bar y(\mu)\right)
   \leq \int v\,\mathrm d\mu;
\]
more explicitly, this is branch convexity between the barycentre of
\(\nu/\beta\), whose weighted radial coordinate is \(r\), and the node
barycentre of the normalized residual measure.  If \(\beta=0\), then
\(r=0\) and the node case gives the same inequality; if \(\beta=1\), it is
ordinary branchwise Jensen.  On an exceptional branch
\(r=s\).  On a full branch \(r\leq s\), so radial nonincrease gives
\[
 v_{i_0}(s,\bar y(\mu))
       \leq v_{i_0}(r,\bar y(\mu))
       \leq\int v\,\mathrm d\mu.
\]
At a full endpoint, apply the finite-root inequality for every sufficiently
large \(s\) and use the exact trace.  This proves
\eqref{eq:p071-finite-exact-jensen-claim}.

Apply the Jensen claim to \(u\) and then take the infimum in
\eqref{eq:p071-finite-exact-moment-envelope}; this proves the first
inequality in \eqref{eq:p071-finite-exact-star-insertion}.  The same claim
shows that every exact comparison function below \(\phi\) is bounded above
by \(E_\phi\), so the envelope is largest.  Monotonicity in the obstacle is
immediate from the formula.  If
\(\|\phi-\psi\|_\infty\leq\delta\), integration over probability measures
gives
\(\|E_\phi-E_\psi\|_\infty\leq\delta\), proving nonexpansion.
\end{proof}
